% Candidate standalone chapter for GAGA source units gaga:no:1 through
% gaga:no:4. This is deliberately the reduced 1956 notion. It does not
% introduce nonreduced complex analytic spaces.

\input{preamble}

% OK, start here.
%
\begin{document}

\title{Complex Analytic Spaces and GAGA}

\maketitle

\phantomsection
\label{section-phantom}

\tableofcontents

\section{Introduction}
\label{section-introduction}

\noindent
This chapter develops the complex-analytic foundations needed to state and
prove the comparison theorems usually called GAGA. The initial sections use
the reduced notion of complex analytic space in
\cite[\S 1, nos. 1--4, pp. 3--7]{GAGA}. This scope is recorded explicitly
because a structure sheaf contained in a sheaf of functions has no
nilpotents. The abstract language of ringed spaces and coherent modules is
that of Sheaves, Definition \ref{sheaves-definition-ringed-space} and
Modules, Section \ref{modules-section-coherent}.

\section{Analytic subsets of complex affine space}
\label{section-analytic-subsets}

\begin{definition}
\label{definition-analytic-subset}
\begin{reference}
\cite[\S 1, no. 1, pp. 3--4]{GAGA}
\end{reference}
Let $n \geq 0$ and endow $\mathbf C^n$ with its usual topology. A subset
$U \subset \mathbf C^n$ is an {\it analytic subset in the sense of Serre}
if, for every $x \in U$, there are an open neighbourhood $W$ of $x$ and
holomorphic functions $f_1, \ldots, f_k$ on $W$ such that
$$
U \cap W = \{z \in W \mid f_1(z)=\ldots=f_k(z)=0\}.
$$

Let $\mathcal C_U$ be the sheaf of germs of all complex-valued functions on
$U$. Let $\mathcal O_{\mathbf C^n}$ be the sheaf of holomorphic functions on
$\mathbf C^n$. The {\it sheaf of holomorphic functions on $U$} is the image
sheaf
$$
\mathcal O_U =
\mathop{\rm Im}\left(
\mathcal O_{\mathbf C^n}|_U \longrightarrow \mathcal C_U
\right).
$$
Thus, if $x \in U$ and $\mathcal I_{U,x}$ is the ideal of germs in
$\mathcal O_{\mathbf C^n,x}$ which vanish on $U$ near $x$, then
$$
\mathcal O_{U,x}=
\mathcal O_{\mathbf C^n,x}/\mathcal I_{U,x}.
$$
\end{definition}

\begin{definition}
\label{definition-holomorphic-map-analytic-subsets}
\begin{reference}
\cite[\S 1, no. 1, pp. 3--4]{GAGA}
\end{reference}
Let $U \subset \mathbf C^r$ and $V \subset \mathbf C^s$ be analytic
subsets. A map $\varphi:U\to V$ is {\it holomorphic} if it is continuous and
composition with $\varphi$ sends every germ in
$\mathcal O_{V,\varphi(x)}$ to a germ in $\mathcal O_{U,x}$ for every
$x\in U$. Equivalently, the $s$ coordinate functions of $\varphi$ are
sections of $\mathcal O_U$ locally on $U$.

A bijective holomorphic map is an {\it analytic isomorphism} if its inverse
is holomorphic.
\end{definition}

\begin{lemma}
\label{lemma-analytic-subset-basic-properties}
Let $U \subset \mathbf C^n$ and $U' \subset \mathbf C^{n'}$ be analytic
subsets in the sense of Definition \ref{definition-analytic-subset}.
\begin{enumerate}
\item The subset $U$ is locally closed and locally compact.
\item The stalk $\mathcal O_{U,x}$ is a reduced local $\mathbf C$-algebra
whose residue homomorphism is evaluation at $x$.
\item The subset $U\times U'$ is analytic in $\mathbf C^{n+n'}$, its
topology is the product topology, and its structure sheaf is the one obtained
from the product analytic charts.
\item Products of holomorphic maps are holomorphic.
\end{enumerate}
\end{lemma}

\begin{proof}
After shrinking around $x\in U$, the set $U$ is the common zero locus of
finitely many continuous functions and hence is closed in an open subset of
$\mathbf C^n$. This proves (1). The displayed quotient in Definition
\ref{definition-analytic-subset} embeds into $\mathcal C_{U,x}$, and a germ
is invertible exactly when its value at $x$ is nonzero. This proves (2).

Near $(x,x')$, concatenate local defining equations for $U$ and $U'$.
This proves analyticity of the product. The remaining assertions in (3) and
(4) follow directly from the induced topology, the coordinate criterion in
Definition \ref{definition-holomorphic-map-analytic-subsets}, and composition
of holomorphic functions.
\end{proof}


\section{Serre complex analytic spaces}
\label{section-serre-analytic-spaces}

\begin{definition}
\label{definition-serre-analytic-space}
\begin{reference}
\cite[\S 1, no. 2, pp. 4--5]{GAGA}
\end{reference}
A {\it complex analytic space in the sense of Serre} is a Hausdorff
topological space $X$ together with a subsheaf of $\mathbf C$-algebras
$$
\mathcal O_X \subset \mathcal C_X
$$
such that $X$ has an open covering $X=\bigcup V_i$ for which every ringed
space $(V_i,\mathcal O_X|_{V_i})$ is analytically isomorphic to an analytic
subset of some $\mathbf C^{n_i}$ endowed with the sheaf of Definition
\ref{definition-analytic-subset}.

A map of such spaces is {\it holomorphic} if it is continuous and pullback
by composition sends germs of holomorphic functions to germs of holomorphic
functions. An isomorphism is a holomorphic map with holomorphic inverse.
\end{definition}

\begin{remark}[Reduced scope]
\label{remark-serre-analytic-spaces-reduced}
Because $\mathcal O_X$ is a subsheaf of a sheaf of functions, every
$\mathcal O_{X,x}$ is reduced. Thus Definition
\ref{definition-serre-analytic-space} is the reduced 1956 notion used by
Serre. It does not include nilpotent analytic thickenings. Any later extension
to nonreduced complex analytic spaces must be introduced separately and must
not be attributed to GAGA, Definition 1.
\end{remark}

\begin{definition}
\label{definition-analytic-subspace}
Let $X$ be a Serre complex analytic space. A subset $Y\subset X$ is an
{\it analytic subspace} if, in every analytic chart $V\to U$, the image of
$Y\cap V$ is an analytic subset of $U$. It has the induced analytic-space
structure. When $Y$ is closed in $X$, we call it a {\it closed analytic
subspace}.
\end{definition}

\begin{lemma}
\label{lemma-serre-analytic-subspaces-products}
\begin{reference}
\cite[\S 1, no. 2, p. 5]{GAGA}
\end{reference}
Let $X$ and $X'$ be Serre complex analytic spaces.
\begin{enumerate}
\item Every analytic subspace of $X$ is locally closed and the induced
structure is independent of the chosen charts.
\item There is a unique Serre complex analytic-space structure on
$X\times X'$ for which products of charts are charts. Its topology is the
product topology.
\item Serre complex analytic spaces and holomorphic maps form a category.
\end{enumerate}
\end{lemma}

\begin{proof}
All three assertions are local on the source spaces. They therefore follow
from Lemma \ref{lemma-analytic-subset-basic-properties}, compatibility of the
restriction maps defining the structure sheaves, and the coordinate
criterion for holomorphic maps.
\end{proof}


\section{Analytic modules and Oka--Cartan coherence}
\label{section-analytic-modules}

\begin{definition}
\label{definition-analytic-module}
\begin{reference}
\cite[\S 1, no. 3, p. 5]{GAGA}
\end{reference}
Let $X$ be a Serre complex analytic space. An {\it analytic module} on $X$
is an $\mathcal O_X$-module. It is {\it coherent} if it is coherent in the
sense of Modules, Definition \ref{modules-definition-coherent}.
\end{definition}

\begin{definition}
\label{definition-vanishing-ideal-analytic-subspace}
Let $i:Y\to X$ be a closed analytic subspace. The {\it vanishing ideal}
$\mathcal J_Y\subset\mathcal O_X$ has stalk
$$
(\mathcal J_Y)_x=
\{f\in\mathcal O_{X,x}\mid f|_Y=0\text{ near }x\}.
$$
There is a canonical isomorphism
$$
\mathcal O_X/\mathcal J_Y \longrightarrow i_*\mathcal O_Y.
$$
\end{definition}

\begin{theorem}[Oka--Cartan coherence]
\label{theorem-oka-cartan-coherence}
\begin{reference}
\cite[\S 1, Proposition 1, pp. 5--6]{GAGA}
\end{reference}
Let $X$ be a Serre complex analytic space.
\begin{enumerate}
\item The $\mathcal O_X$-module $\mathcal O_X$ is coherent.
\item If $Y\subset X$ is a closed analytic subspace, then its vanishing
ideal $\mathcal J_Y$ is a coherent $\mathcal O_X$-module.
\end{enumerate}
\end{theorem}

\begin{proof}
The analytic input is the Oka--Cartan theorem that, for an open subset
$\Omega\subset\mathbf C^n$, the sheaf $\mathcal O_\Omega$ and the vanishing
ideal of every closed analytic subset of $\Omega$ are coherent.

The assertions are local on $X$. Choose a chart and shrink it so that
$X=Z(\mathcal I_X)$ is a closed analytic subset of an open
$\Omega\subset\mathbf C^n$. If $Y\subset X$ is closed analytic, shrink again
and write $Y=Z(\mathcal I_Y)$ with
$\mathcal I_X\subset\mathcal I_Y\subset\mathcal O_\Omega$. Then
$$
\mathcal O_X=\mathcal O_\Omega/\mathcal I_X,
\qquad
\mathcal J_Y=\mathcal I_Y/\mathcal I_X.
$$
The Oka--Cartan input makes $\mathcal O_\Omega$, $\mathcal I_X$, and
$\mathcal I_Y$ coherent. Modules, Lemmas
\ref{modules-lemma-coherent-abelian} and
\ref{modules-lemma-coherent-change-rings} now give (1) and (2).
\end{proof}

\begin{remark}[Examples]
\label{remark-examples-coherent-analytic-modules}
The sheaf of holomorphic sections of a holomorphic vector bundle is finite
locally free and hence coherent by Theorem
\ref{theorem-oka-cartan-coherence}; compare Cohomology, Lemma
\ref{cohomology-lemma-holomorphic-vector-bundles-locally-free-analytic-modules}.
Serre also records the sheaves of germs of automorphic functions from the
Cartan seminar as coherent examples; this cited analytic construction is not
needed for the foundation statements above.
\end{remark}


\section{Local analytic algebras}
\label{section-local-analytic-algebras}

\begin{lemma}
\label{lemma-local-analytic-algebra}
\begin{reference}
\cite[\S 1, no. 4, p. 6]{GAGA}
\end{reference}
Let $X$ be a Serre complex analytic space and $x\in X$. There are an
integer $n\geq 0$ and a radical ideal
$\mathfrak a\subset\mathbf C\{z_1,\ldots,z_n\}$ such that
$$
\mathcal O_{X,x}\cong
\mathbf C\{z_1,\ldots,z_n\}/\mathfrak a
$$
as local $\mathbf C$-algebras. Its maximal ideal consists of germs vanishing
at $x$, its residue field is $\mathbf C$, and it is reduced and Noetherian.
\end{lemma}

\begin{proof}
Choose an analytic chart at $x$ and translate its image point to the origin.
Definition \ref{definition-analytic-subset} gives the displayed quotient.
The ideal is radical because it is the ideal of germs vanishing on a set.
The residue and reducedness statements also follow from Lemma
\ref{lemma-analytic-subset-basic-properties}. The remaining analytic input is
the Weierstrass theorem that
$\mathbf C\{z_1,\ldots,z_n\}$ is Noetherian; the quotient is then Noetherian
by Algebra, Lemma \ref{algebra-lemma-Noetherian-permanence}.
\end{proof}

\begin{lemma}
\label{lemma-analytic-germ-local-algebra}
\begin{reference}
\cite[\S 1, no. 4, p. 6]{GAGA}
\end{reference}
The local analytic $\mathbf C$-algebra $\mathcal O_{X,x}$ determines the
germ of a Serre complex analytic space $X$ at $x$. More precisely, an
isomorphism of local analytic $\mathbf C$-algebras
$\mathcal O_{X,x}\cong\mathcal O_{Y,y}$ is induced contravariantly by an
isomorphism of analytic germs $(X,x)\cong(Y,y)$.
\end{lemma}

\begin{proof}
Present both local algebras as quotients of convergent power-series rings.
Images of the coordinate germs under a local homomorphism are convergent
germs and hence define a holomorphic map germ. The relations force its image
to lie in the required analytic germ. Applying the same construction to the
inverse algebra homomorphism gives inverse holomorphic germs.
\end{proof}

\begin{lemma}
\label{lemma-simple-point-regular-local}
\begin{reference}
\cite[\S 1, no. 4, p. 6]{GAGA}
\end{reference}
Let $X$ be a Serre complex analytic space and $x\in X$. The following are
equivalent for $n\geq 0$:
\begin{enumerate}
\item the germ $(X,x)$ is isomorphic to $(\mathbf C^n,0)$;
\item $\mathcal O_{X,x}\cong\mathbf C\{z_1,\ldots,z_n\}$;
\item $\mathcal O_{X,x}$ is a regular local ring of dimension $n$ in the
sense of Algebra, Definition \ref{algebra-definition-regular-local}.
\end{enumerate}
Such an $x$ is a {\it simple point of dimension $n$}. If every point of $X$
is simple, then $X$ is an {\it analytic variety}.
\end{lemma}

\begin{proof}
The equivalence of (1) and (2) follows from Lemma
\ref{lemma-analytic-germ-local-algebra}. The implication (2) $\Rightarrow$
(3) is the standard dimension computation for convergent power-series
rings. The converse is the regular analytic local-algebra theorem, proved
from the analytic inverse-function theorem (equivalently, by choosing a
basis of $\mathfrak m_x/\mathfrak m_x^2$ and applying the analytic implicit
function theorem). This analytic theorem is an explicit external dependency
of the present foundation block.
\end{proof}

\begin{lemma}
\label{lemma-components-dimension-analytic-space}
\begin{reference}
\cite[\S 1, no. 4, pp. 6--7]{GAGA}
\end{reference}
Let $X$ be a Serre complex analytic space and $x\in X$.
\begin{enumerate}
\item The irreducible components $X_i$ of the germ $(X,x)$ correspond to
the finitely many minimal primes $\mathfrak p_i$ of $\mathcal O_{X,x}$, and
$$
0=\bigcap_i\mathfrak p_i,
\qquad
\mathcal O_{X_i,x}=\mathcal O_{X,x}/\mathfrak p_i.
$$
\item The analytic dimension of $X$ at $x$ (one half of its local topological
dimension) is the supremum of the dimensions of the $X_i$ and equals the
Krull dimension of $\mathcal O_{X,x}$.
\end{enumerate}
\end{lemma}

\begin{proof}
Reducedness and Noetherianity are given by Lemma
\ref{lemma-local-analytic-algebra}. The algebraic assertion about the
finiteness and intersection of minimal primes follows from Algebra, Lemmas
\ref{algebra-lemma-Noetherian-irreducible-components} and
\ref{algebra-lemma-Zariski-topology}. The correspondence with analytic
components uses the local analytic Nullstellensatz. For (2), the analytic
normalization/parameter theorem says that an irreducible analytic local
algebra of analytic dimension $r$ is finite over
$\mathbf C\{z_1,\ldots,z_r\}$. Dimension invariance under finite integral
extensions then gives Krull dimension $r$. The analytic Nullstellensatz and
local parameter theorem are explicit external dependencies; no algebraic
Stacks result alone supplies them.
\end{proof}




% Candidate continuation fragment for gaga.tex covering GAGA source units
% gaga:no:5 through gaga:lemma:3. Insert this fragment immediately before
% the unique terminal \input{chapters} in the foundation candidate.
%
% Convention: throughout this fragment an algebraic variety over C means a
% reduced separated scheme of finite type over C, not necessarily irreducible.

\section{Analytification of algebraic varieties}
\label{section-analytification}

\noindent
Throughout this section an {\it algebraic variety over $\mathbf C$} means a
reduced separated scheme of finite type over $\mathbf C$, not necessarily
irreducible. Its classical points are its closed points, equivalently its
$\mathbf C$-points. This convention is the scheme-theoretic translation of
the convention in \cite[\S 2, no. 5]{GAGA} and is compatible with the reduced
analytic spaces of Section \ref{section-serre-analytic-spaces}.

\begin{lemma}
\label{lemma-algebraic-charts-analytic}
Let $U \subset \mathbf C^n$ and $U' \subset \mathbf C^{n'}$ be Zariski locally
closed reduced subsets.
\begin{enumerate}
\item The Zariski topology on $\mathbf C^n$ is coarser than its usual topology.
\item The set $U$, with its reduced induced structure, is a Serre complex
analytic space.
\item A regular map $f : U \to U'$ is holomorphic.
\item A biregular isomorphism $U \to U'$ is an analytic isomorphism.
\end{enumerate}
\end{lemma}

\begin{proof}
A Zariski closed subset is the common zero set of polynomials, and polynomials
are continuous and holomorphic. This proves (1) and proves (2) for a closed
subset; a locally closed subset is open in its closure, and restriction gives
the general case. A regular function on a Zariski open subset is locally a
quotient $g/h$ of polynomials with $h$ nonvanishing. Hence it is holomorphic.
Applying this coordinate by coordinate proves (3). Applying (3) to $f$ and to
$f^{-1}$ proves (4).
\end{proof}

\begin{proposition}
\label{proposition-analytification}
Let $X$ be an algebraic variety over $\mathbf C$. There is a unique Serre
complex analytic space $X^{an}$ with the following property: for every
algebraic chart
$$
\varphi : V \longrightarrow U \subset \mathbf C^n
$$
from a Zariski open $V \subset X$ onto a Zariski locally closed reduced
subset, $V$ is open in $X^{an}$ and $\varphi$ is an analytic isomorphism.

The construction is functorial. Every morphism $f : X \to Y$ induces a
holomorphic map $f^{an} : X^{an} \to Y^{an}$, and there is a canonical
morphism of locally ringed spaces
$$
a_X : (X^{an}, \mathcal O_{X^{an}}) \longrightarrow
(X, \mathcal O_X)
$$
whose map on points sends a classical point to the corresponding closed
scheme point and whose map on structure sheaves sends a regular function to
the same function viewed as holomorphic.
\end{proposition}

\begin{proof}
Cover $X$ by finitely many algebraic charts $V_i \to U_i$. Transport the
analytic structure of $U_i$ to $V_i$. On $V_i \cap V_j$ the transition maps
and their inverses are holomorphic by
Lemma \ref{lemma-algebraic-charts-analytic}; the overlap is open for the usual
topology by the same lemma. The cocycle condition is inherited from the
algebraic charts. The sheaves and the locally ringed spaces therefore glue by
Sheaves, Lemma \ref{sheaves-lemma-glue-sheaves-structures}.

It remains to check the Hausdorff axiom, rather than hide it in the word
``glue''. The inverse image of the diagonal of the glued point set in every
$V_i \times V_j$ is the graph of the algebraic identification on the overlap.
Since $X$ is separated, this graph is Zariski closed; it is closed in the usual
topology by Lemma \ref{lemma-algebraic-charts-analytic}. These products form an
open cover, so the diagonal is closed and the glued space is Hausdorff.
Uniqueness follows because the charts cover $X$ and determine both its topology
and its structure sheaf.

For a morphism $f : X \to Y$, its restrictions in algebraic charts are
holomorphic by Lemma \ref{lemma-algebraic-charts-analytic}; they glue to
$f^{an}$, compatibly with identities and composition. The usual topology is
finer than the Zariski topology, so the point map underlying $a_X$ is
continuous. Regular functions are holomorphic in charts, and the induced maps
on stalks are local. This constructs the asserted morphism of locally ringed
spaces.
\end{proof}

\begin{lemma}
\label{lemma-analytification-properties}
Let $X$ and $Y$ be algebraic varieties over $\mathbf C$.
\begin{enumerate}
\item The space $X^{an}$ is locally compact, second countable, and countable
at infinity.
\item There is a canonical analytic isomorphism
$(X \times_{\mathbf C} Y)^{an} \cong X^{an} \times Y^{an}$.
\item If $Z \subset X$ is a Zariski locally closed reduced subvariety, then
$Z^{an}$ is the reduced analytic subspace induced on $Z(\mathbf C)$ by
$X^{an}$.
\end{enumerate}
\end{lemma}

\begin{proof}
A Zariski locally closed subset of $\mathbf C^n$ is locally compact, second
countable, and a countable union of compact subsets. The same properties pass
to a finite union of open chart domains, proving (1). Products of algebraic
charts are algebraic charts and the usual topology on a product of subsets of
complex affine spaces is the product topology; uniqueness in
Proposition \ref{proposition-analytification} proves (2). Restricting charts of
$X$ to $Z$ and using Lemma \ref{lemma-algebraic-charts-analytic} proves (3).
\end{proof}

\section{Local algebra of analytification}
\label{section-local-comparison}

\noindent
Let $x \in X(\mathbf C)$. We write
$$
\theta_x : \mathcal O_{X,x} \longrightarrow \mathcal O_{X^{an},x}
$$
for the local homomorphism induced by $a_X$. Both rings are Noetherian local
rings with residue field $\mathbf C$: this is standard for the algebraic ring,
and for the analytic ring it is Lemma \ref{lemma-local-analytic-algebra}.

\begin{proposition}
\label{proposition-analytification-closed-ideal}
Let $Z \subset X$ be a Zariski locally closed reduced subvariety and let
$x \in Z(\mathbf C)$. After replacing $X$ by a Zariski neighbourhood of $x$,
view $Z$ as closed. If $\mathcal J_{Z,x} \subset \mathcal O_{X,x}$ is its
algebraic ideal and $\mathcal I_{Z^{an},x} \subset
\mathcal O_{X^{an},x}$ is the ideal of holomorphic germs vanishing on
$Z^{an}$, then
$$
\mathcal I_{Z^{an},x} =
\mathcal J_{Z,x}\mathcal O_{X^{an},x}.
$$
Consequently
$$
\mathcal O_{Z^{an},x} \cong
\mathcal O_{X^{an},x}/
\mathcal J_{Z,x}\mathcal O_{X^{an},x}.
$$
\end{proposition}

\begin{proof}
First take $X=\mathbf A^n_{\mathbf C}$ and translate $x$ to the origin. Put
$$
R=\mathbf C[z_1,\ldots,z_n]_{(z_1,\ldots,z_n)},
\qquad
H=\mathbf C\{z_1,\ldots,z_n\}.
$$
Let $J \subset R$ be the radical ideal of $Z$. The local analytic
Nullstellensatz says that the ideal of holomorphic germs vanishing on $Z$ is
$\sqrt{JH}$. We show that $JH$ is already radical.

The Taylor homomorphisms identify the maximal-adic completions of both $R$ and
$H$ with $\mathbf C[[z_1,\ldots,z_n]]$. Completion is exact on finite modules
by Algebra, Lemma \ref{algebra-lemma-completion-tensor}; hence the completion
of $H/JH$ is
$$
\mathbf C[[z_1,\ldots,z_n]]/
J\mathbf C[[z_1,\ldots,z_n]],
$$
which is also the completion of $R/J$. The field $\mathbf C$ is excellent by
More on Algebra, Proposition
\ref{more-algebra-proposition-ubiquity-excellent}, and $R/J$, being a
localization of a finite type $\mathbf C$-algebra, is excellent by More on
Algebra, Lemma \ref{more-algebra-lemma-finite-type-over-excellent}. It is
therefore Nagata by More on Algebra, Lemma
\ref{more-algebra-lemma-quasi-excellent-nagata}. Its
completion is reduced by More on Algebra, Lemma
\ref{more-algebra-lemma-completion-reduced}. Thus the completion of $H/JH$ is
reduced. The ring $H/JH$ injects into its completion by Krull's intersection
theorem, Algebra, Lemma
\ref{algebra-lemma-intersect-powers-ideal-module-zero}; hence $H/JH$ is
reduced and $JH=\sqrt{JH}$. The analytic Nullstellensatz now gives the desired
equality in affine space.

For general $X$, work in an algebraic chart which realizes $X$ and $Z$ as
reduced locally closed subsets of one affine space. Apply the affine result to
their ambient ideals and pass to the quotient. This gives both displayed
formulas.
\end{proof}

\begin{theorem}
\label{theorem-analytification-completed-local-rings}
For every algebraic variety $X$ over $\mathbf C$ and every
$x \in X(\mathbf C)$, the local homomorphism $\theta_x$ induces an isomorphism
of completed local rings
$$
\widehat{\theta}_x :
\widehat{\mathcal O}_{X,x} \xrightarrow{\ \sim\ }
\widehat{\mathcal O}_{X^{an},x}.
$$
In particular, $\theta_x$ is injective.
\end{theorem}

\begin{proof}
For affine space at the origin the map is
$$
\mathbf C[z_1,\ldots,z_n]_{(z_1,\ldots,z_n)}
\longrightarrow \mathbf C\{z_1,\ldots,z_n\},
$$
and both maximal-adic completions are canonically
$\mathbf C[[z_1,\ldots,z_n]]$ by Taylor expansion and truncation. For a
locally closed reduced $X$ in affine space, Proposition
\ref{proposition-analytification-closed-ideal} identifies its analytic local
ring with the quotient by the extended algebraic ideal. Exactness of
completion on finite modules, Algebra, Lemma
\ref{algebra-lemma-completion-tensor}, identifies both completions with the
same quotient of the formal power-series ring. The assertion is local, so
charts prove the general case. Finally, $\mathcal O_{X,x}$ injects into its
completion by Algebra, Lemma
\ref{algebra-lemma-intersect-powers-ideal-module-zero}, which proves the last
sentence.
\end{proof}

\begin{theorem}
\label{theorem-analytification-local-faithfully-flat}
The local homomorphism
$$
\theta_x : \mathcal O_{X,x} \longrightarrow \mathcal O_{X^{an},x}
$$
is faithfully flat. In Serre's terminology,
$(\mathcal O_{X,x},\mathcal O_{X^{an},x})$ is a flat pair: the quotient
$\mathcal O_{X^{an},x}/\mathcal O_{X,x}$ is flat over $\mathcal O_{X,x}$,
and for every ideal $I \subset \mathcal O_{X,x}$ one has
$$
I\mathcal O_{X^{an},x}\cap\mathcal O_{X,x}=I.
$$
\end{theorem}

\begin{proof}
The map on completions is an isomorphism by
Theorem \ref{theorem-analytification-completed-local-rings}, hence is flat.
More on Algebra, Lemma \ref{more-algebra-lemma-flat-completion} implies that
$\theta_x$ is flat, and Algebra, Lemma
\ref{algebra-lemma-local-flat-ff} makes this local flat map faithfully flat.
The contraction formula is Algebra, Lemma
\ref{algebra-lemma-faithfully-flat-universally-injective}. Universal
injectivity after tensoring the exact sequence
$0\to\mathcal O_{X,x}\to\mathcal O_{X^{an},x}$ shows that the displayed
quotient is flat, which is precisely the flat-pair assertion.
\end{proof}

\begin{lemma}
\label{lemma-analytification-local-dimension}
For every $x \in X(\mathbf C)$,
$$
\dim(\mathcal O_{X,x})=\dim(\mathcal O_{X^{an},x}).
$$
\end{lemma}

\begin{proof}
Each ring has the same dimension as its completion by More on Algebra, Lemma
\ref{more-algebra-lemma-completion-dimension}, and the two completions are
isomorphic by Theorem \ref{theorem-analytification-completed-local-rings}.
\end{proof}

\begin{lemma}
\label{lemma-analytification-dimension}
If $X$ is irreducible of dimension $r$, then $X^{an}$ has analytic dimension
$r$ at every point.
\end{lemma}

\begin{proof}
For a closed point $x$, Varieties, Lemma
\ref{varieties-lemma-dimension-locally-algebraic}, parts
\ref{varieties-item-dimension-at-closed-point} and
\ref{varieties-item-dimension-irreducible}, gives
$$
\dim(\mathcal O_{X,x})=r.
$$
Apply
Lemma \ref{lemma-analytification-local-dimension}, followed by
Lemma \ref{lemma-components-dimension-analytic-space}, which identifies
analytic dimension with the Krull dimension of the analytic local ring.
\end{proof}

\section{The classical and Zariski topologies}
\label{section-classical-zariski-topologies}

\begin{lemma}
\label{lemma-zariski-dense-open-analytically-dense}
Let $U \subset X$ be Zariski open and Zariski dense. Then $U^{an}$ is dense
in $X^{an}$.
\end{lemma}

\begin{proof}
Put $Z=X\setminus U$ with its reduced induced structure. If a point
$x \in X^{an}$ were not in the closure of $U^{an}$, then $Z^{an}$ would
contain an analytic neighbourhood of $x$. Its local vanishing ideal at $x$
would therefore be zero. Proposition
\ref{proposition-analytification-closed-ideal} and the injectivity in
Theorem \ref{theorem-analytification-completed-local-rings} would give
$\mathcal J_{Z,x}=0$. Thus $Z$ would agree with $X$ on a nonempty Zariski
neighbourhood of $x$, contradicting the Zariski density of $U$.
\end{proof}

\begin{proposition}
\label{proposition-proper-if-and-only-if-compact-analytification}
An algebraic variety $X$ over $\mathbf C$ is proper over $\mathbf C$ if and
only if the topological space $X^{an}$ is compact.
\end{proposition}

\begin{proof}
Apply Chow's lemma in the precise form of Limits, Lemma
\ref{limits-lemma-chow-finite-type}. After replacing the source by its
reduction as allowed by the following remark there, we obtain a proper
surjective morphism $\pi : U \to X$ and an immersion
$U \to \mathbf P^n_{\mathbf C}$. Let $Y$ be the reduced closure of $U$ in
$\mathbf P^n_{\mathbf C}$. Then $Y$ is projective and $U$ is Zariski open
and dense in $Y$. Moreover the graph of $\pi$ is closed in $X\times Y$:
over $X$ it is the graph of a morphism from the proper $X$-scheme $U$ to the
separated $X$-scheme $X\times Y$.

If $X$ is proper, then $U$ is proper over $\mathbf C$. Its immersion into
projective space is proper and hence closed, so $U=Y$. The space
$Y^{an}$ is a closed subset of the compact Hausdorff space
$\mathbf P^n(\mathbf C)$, hence compact. The continuous surjection
$\pi^{an}:Y^{an}\to X^{an}$ shows that $X^{an}$ is compact.

Conversely, assume $X^{an}$ compact. Then $X^{an}\times Y^{an}$ is compact,
and the analytification of the closed graph is closed in this product by
Lemma \ref{lemma-analytification-properties}. Its projection to $Y^{an}$ is
$U^{an}$, hence $U^{an}$ is compact and therefore closed in the Hausdorff
space $Y^{an}$. It is dense by
Lemma \ref{lemma-zariski-dense-open-analytically-dense}, so $U^{an}=Y^{an}$.
Thus $U=Y$ because a nonempty closed subset of a finite type
$\mathbf C$-scheme has a closed $\mathbf C$-point. Now the projective variety
$Y$ is proper and maps surjectively to $X$; Morphisms, Lemma
\ref{morphisms-lemma-image-proper-is-proper} implies that $X$ is proper.
\end{proof}

\begin{lemma}
\label{lemma-dense-image-contains-open}
Let $f : X \to Y$ be a morphism of algebraic varieties over $\mathbf C$. If
$f(X)$ is Zariski dense in $Y$, then $f(X)$ contains a Zariski open dense
subset of $Y$.
\end{lemma}

\begin{proof}
The image is constructible by Morphisms, Lemma
\ref{morphisms-lemma-chevalley}. A dense constructible subset contains a dense
open subset by Topology, Lemma
\ref{topology-lemma-dense-in-constructible}.
\end{proof}

\begin{proposition}
\label{proposition-image-closure}
Let $f : X \to Y$ be a morphism of algebraic varieties. The closure of
$f(X(\mathbf C))$ in $Y^{an}$ is the set of classical points of the Zariski
closure of $f(X)$ in $Y$.
\end{proposition}

\begin{proof}
Let $T$ be the reduced Zariski closure of $f(X)$. Applying
Lemma \ref{lemma-dense-image-contains-open} to $X\to T$ gives a Zariski open
dense $V\subset T$ contained in $f(X)$. By
Lemma \ref{lemma-zariski-dense-open-analytically-dense}, $V^{an}$ is dense in
$T^{an}$, so $T^{an}$ is contained in the analytic closure of $f(X)$. The
opposite inclusion holds because $T^{an}$ is closed in $Y^{an}$ by
Lemma \ref{lemma-analytification-properties}.
\end{proof}

\section{An analytic criterion for regularity}
\label{section-analytic-regularity}

\begin{proposition}
\label{proposition-algebraic-graph-holomorphic-regular}
Let $X$ and $Y$ be algebraic varieties and let
$f : X^{an}\to Y^{an}$ be holomorphic. If its graph is the classical point
set of a Zariski locally closed reduced subvariety
$T\subset X\times_{\mathbf C}Y$, then $f$ is induced by a morphism
$X\to Y$.
\end{proposition}

\begin{proof}
The first projection $p:T\to X$ is a morphism and is bijective on classical
points. Its inverse $x\mapsto(x,f(x))$ is holomorphic. By
Lemma \ref{lemma-analytification-properties}, $p^{an}$ is therefore an
analytic isomorphism. Proposition
\ref{proposition-analytic-isomorphism-algebraic} below shows that $p$ is an
isomorphism. Hence $f=\operatorname{pr}_Y\circ p^{-1}$ is regular.
\end{proof}

\begin{proposition}
\label{proposition-analytic-isomorphism-algebraic}
Let $p:T\to X$ be a morphism of algebraic varieties which is bijective on
classical points. If $p^{an}:T^{an}\to X^{an}$ is an analytic isomorphism,
then $p$ is an isomorphism.
\end{proposition}

\begin{proof}
First $p$ is a homeomorphism for the Zariski topology on classical points.
Indeed, if $F\subset T$ is Zariski closed, then $F^{an}$ is closed in
$T^{an}$. Its image under the analytic homeomorphism is closed in $X^{an}$.
By Proposition \ref{proposition-image-closure}, applied to $F\to X$, this
image is Zariski closed.

Both schemes are Jacobson by Varieties, Lemma
\ref{varieties-lemma-locally-finite-type-Jacobson}. Hence their closed subsets
are determined by their closed points. Passing from irreducible closed subsets
to their unique generic points (Schemes, Lemma
\ref{schemes-lemma-scheme-sober}) extends the preceding homeomorphism uniquely
to a homeomorphism of the full underlying scheme spaces.

Fix $t\in T(\mathbf C)$ and put $x=p(t)$. The Zariski homeomorphism makes
$$
A=\mathcal O_{X,x}\longrightarrow A'=\mathcal O_{T,t}
$$
injective. The analytic isomorphism identifies
$B=\mathcal O_{X^{an},x}$ with $\mathcal O_{T^{an},t}$. By
Theorem \ref{theorem-analytification-local-faithfully-flat}, both $A$ and
$A'$ embed in this common ring $B$.

The irreducible components through $x$ correspond to those through $t$.
Write their minimal primes as $\mathfrak p_i\subset A$ and
$\mathfrak p'_i\subset A'$, with
$\mathfrak p'_i\cap A=\mathfrak p_i$. Put
$$
K_i=\operatorname{Frac}(A/\mathfrak p_i),
\qquad
K'_i=\operatorname{Frac}(A'/\mathfrak p'_i).
$$
The corresponding irreducible varieties are homeomorphic and have the same
dimension, so $K'_i/K_i$ is a finite extension. To see that its degree is one,
use Morphisms, Lemma \ref{morphisms-lemma-finite-degree} and generic flatness,
Morphisms, Proposition \ref{morphisms-proposition-generic-flatness}, to shrink
the target to a nonempty open over which the component map is finite locally
free. Since the ground field has characteristic zero, $K'_i/K_i$ is separable;
after shrinking once more the map is finite etale of degree
$[K'_i:K_i]$. Every closed fibre then consists of that many reduced
$\mathbf C$-points. Bijectivity of $p$ forces $[K'_i:K_i]=1$.

Let $S$ be the complement in $A$ of the union of its minimal primes and define
$S'$ similarly for $A'$. Lemma \ref{lemma-total-fractions-reduced} below gives
$$
A_S=\prod_iK_i=\prod_iK'_i=A'_{S'}.
$$
If $h\in A'$, write $h=g/s$ in this common total ring of fractions with
$g\in A$ and $s\in S$. Thus $g=sh\in sB$. Faithful flatness and Algebra,
Lemma \ref{algebra-lemma-faithfully-flat-universally-injective}, give
$sB\cap A=sA$, so $g=sa$ for some $a\in A$. Hence $s(h-a)=0$. The element
$s$ lies in no minimal prime of the reduced ring $A'$, and is therefore a
nonzerodivisor by Algebra, Lemma
\ref{algebra-lemma-reduced-ring-sub-product-fields}. Thus $h=a$ and $A'=A$.

We have proved that the local maps are isomorphisms at every closed point.
Thus $p$ is etale at every closed point by Etale, Lemma
\ref{etale-lemma-characterize-etale-Noetherian}. The etale locus is open; its
closed complement, if nonempty, would contain a closed point because $T$ is
Jacobson. Hence $p$ is etale everywhere.

Finally, repeat the finite-degree argument above for the restriction of $p$ to
the reduced closure of an arbitrary point of $T$. The Zariski homeomorphism
identifies this closure with the closure of its image, and the same argument
shows that the induced extension of their function fields, namely the residue
field extension at the two generic points, has degree one. Thus $p$ is
injective on points with trivial residue extensions, hence radicial by
Morphisms, Lemma \ref{morphisms-lemma-universally-injective}. An etale
radicial morphism is an open immersion by Etale, Theorem
\ref{etale-theorem-etale-radicial-open}. Since $p$ is also surjective, it is an
isomorphism.
\end{proof}

\begin{lemma}
\label{lemma-total-fractions-reduced}
Let $A$ be a reduced ring with finitely many minimal primes
$\mathfrak p_1,\ldots,\mathfrak p_r$. Let $S$ be the set of elements lying in
none of the $\mathfrak p_i$ and put
$K_i=\operatorname{Frac}(A/\mathfrak p_i)$. Then
$$
A_S=\prod_{i=1}^rK_i.
$$
\end{lemma}

\begin{proof}
Algebra, Lemma \ref{algebra-lemma-reduced-ring-sub-product-fields} says that
$S$ is exactly the set of nonzerodivisors and that
$A_{\mathfrak p_i}=K_i$. Algebra, Lemma
\ref{algebra-lemma-total-ring-fractions-no-embedded-points} then identifies
the total ring of fractions $A_S$ with the displayed product.
\end{proof}

% External analytic input ledger for this continuation:
% (1) the local analytic Nullstellensatz, used exactly once in the proof of
%     Proposition \ref{proposition-analytification-closed-ideal};
% (2) Taylor expansion and the completion identity
%     \widehat{\mathbf C\{z_1,...,z_n\}}=\mathbf C[[z_1,...,z_n]], used in
%     the two local-comparison proofs;
% (3) compactness and Hausdorffness of complex projective space, used in the
%     proof of Proposition
%     \ref{proposition-proper-if-and-only-if-compact-analytification}.
% Oka--Cartan coherence and Weierstrass Noetherianity were already stated as
% explicit dependencies in the foundation block; no properness-comparison,
% identity, normalization, or unmentioned analytic theorem is invoked here.

% Continuation fragment for GAGA source nos. 9--12. Insert after the
% analytification fragment and before the proof fragment for nos. 13--17.

\section{Analytification of modules}
\label{section-analytification-modules}

\begin{definition}
\label{definition-analytification-module}
\begin{reference}
\cite[\S 3, no. 9, Definition 2]{GAGA}
\end{reference}
Let $X$ be an algebraic variety over $\mathbf C$ and let
$a_X:X^{\mathrm{an}}\to X$ be the canonical morphism of locally ringed
spaces of Proposition \ref{proposition-analytification}. For an
$\mathcal O_X$-module $\mathcal F$, its {it analytification} is
$$
\mathcal F^{\mathrm{an}}=a_X^*\mathcal F
=\mathcal O_{X^{\mathrm{an}}}
\otimes_{a_X^{-1}\mathcal O_X}a_X^{-1}\mathcal F.
$$
There is a canonical map of $a_X^{-1}\mathcal O_X$-modules
$$
\eta_{\mathcal F}:a_X^{-1}\mathcal F\longrightarrow
\mathcal F^{\mathrm{an}},\qquad s\longmapsto1\otimes s.
$$
Every $\mathcal O_X$-linear map $\mathcal F\to\mathcal G$ induces an
$\mathcal O_{X^{\mathrm{an}}}$-linear map
$\mathcal F^{\mathrm{an}}\to\mathcal G^{\mathrm{an}}$, functorially, and
$\mathcal O_X^{\mathrm{an}}=\mathcal O_{X^{\mathrm{an}}}$ canonically.
\end{definition}

\begin{lemma}
\label{lemma-analytification-module-stalks}
Let $x\in X(\mathbf C)=X^{\mathrm{an}}$. There are canonical identifications
$$
(\mathcal F^{\mathrm{an}})_x
=\mathcal F_x\otimes_{\mathcal O_{X,x}}
\mathcal O_{X^{\mathrm{an}},x}
$$
under which $(\eta_{\mathcal F})_x$ is the map $m\mapsto1\otimes m$.
\end{lemma}

\begin{proof}
Stalks commute with inverse image and tensor product. Since $x$ is a closed
point, $(a_X^{-1}\mathcal F)_x=\mathcal F_x$. The displayed identification
and the description of $\eta_{\mathcal F}$ follow directly from Definition
\ref{definition-analytification-module}.
\end{proof}

\begin{lemma}
\label{lemma-analytification-exact-faithful-coherent}
\begin{reference}
\cite[\S 3, no. 9, Proposition 10]{GAGA}
\end{reference}
Let $X$ be an algebraic variety over $\mathbf C$.
\begin{enumerate}
\item The functor $\mathcal F\mapsto\mathcal F^{\mathrm{an}}$ on
$\mathcal O_X$-modules is exact.
\item The map $\eta_{\mathcal F}:a_X^{-1}\mathcal F\to
\mathcal F^{\mathrm{an}}$ is injective for every $\mathcal F$.
\item If $\mathcal F$ is coherent, then $\mathcal F^{\mathrm{an}}$ is a
coherent analytic module.
\end{enumerate}
\end{lemma}

\begin{proof}
The stalk formula of Lemma \ref{lemma-analytification-module-stalks} and
faithful flatness in Theorem
\ref{theorem-analytification-local-faithfully-flat} prove (1); compare
Modules, Lemma \ref{modules-lemma-pullback-flat}. The same stalk formula and
Algebra, Lemma
\ref{algebra-lemma-faithfully-flat-universally-injective} prove (2).

For (3), a coherent module on the locally Noetherian scheme $X$ is locally
finitely presented. Pullback preserves finite presentations. The analytic
structure sheaf is coherent by Theorem
\ref{theorem-oka-cartan-coherence}; hence a finitely presented analytic
module is coherent by Modules, Lemma
\ref{modules-lemma-coherent-structure-sheaf}.
\end{proof}

\begin{remark}[Analytification of an ideal]
\label{remark-analytification-ideal-sheaf}
If $\mathcal J\subset\mathcal O_X$ is a quasi-coherent ideal, exactness
identifies $\mathcal J^{\mathrm{an}}$ with the ideal of
$\mathcal O_{X^{\mathrm{an}}}$ generated by the image of
$a_X^{-1}\mathcal J$. For the ideal of a reduced closed subvariety this is
also Proposition \ref{proposition-analytification-closed-ideal}.
\end{remark}


\section{Closed immersions and analytification}
\label{section-analytification-closed-immersions}

\begin{lemma}
\label{lemma-analytification-pushforward-closed-immersion}
\begin{reference}
\cite[\S 3, no. 10, Proposition 11]{GAGA}
\end{reference}
Let $i:Z\to X$ be a closed immersion of algebraic varieties over
$\mathbf C$, and let $\mathcal F$ be a coherent $\mathcal O_Z$-module. There
is a canonical isomorphism
$$
(i_*\mathcal F)^{\mathrm{an}}
\longrightarrow i^{\mathrm{an}}_*(\mathcal F^{\mathrm{an}}).
$$
\end{lemma}

\begin{proof}
The commutative analytification square and adjunction give the canonical
map. Both sides vanish away from $Z^{\mathrm{an}}$. At
$z\in Z(\mathbf C)$ the source stalk is
$$
\mathcal F_z\otimes_{\mathcal O_{X,z}}
\mathcal O_{X^{\mathrm{an}},z},
$$
whereas the target stalk is
$$
\mathcal F_z\otimes_{\mathcal O_{Z,z}}
\mathcal O_{Z^{\mathrm{an}},z}.
$$
Proposition \ref{proposition-analytification-closed-ideal} gives
$$
\mathcal O_{Z^{\mathrm{an}},z}
=\mathcal O_{Z,z}\otimes_{\mathcal O_{X,z}}
\mathcal O_{X^{\mathrm{an}},z}.
$$
Associativity of tensor product identifies the two stalks, and the
identification is the stalk of the canonical map.
\end{proof}


\section{The cohomology comparison map}
\label{section-analytification-cohomology-map}

\begin{definition}[Cohomology comparison map]
\label{construction-analytification-cohomology-map}
\begin{reference}
\cite[\S 3, no. 11]{GAGA}
\end{reference}
Let $\mathcal F$ be an $\mathcal O_X$-module. Exactness of $a_X^*$ and the
derived adjunction unit give
$$
\mathcal F\longrightarrow Ra_{X,*}a_X^*\mathcal F.
$$
After applying derived global sections and using
$$
R\Gamma(X,Ra_{X,*}\mathcal G)
=R\Gamma(X^{\mathrm{an}},\mathcal G),
$$
we obtain canonical maps
$$
\epsilon^q_{\mathcal F}:
H^q(X,\mathcal F)\longrightarrow
H^q(X^{\mathrm{an}},\mathcal F^{\mathrm{an}}),
\qquad q\geq0.
$$
For $q=0$, this sends a section $s$ to $1\otimes s$. Thus it agrees with
Serre's map obtained from Zariski covers and their analytifications.
\end{definition}

\begin{lemma}
\label{lemma-analytification-cohomology-map-functorial}
The maps $\epsilon^q_{\mathcal F}$ are natural in $\mathcal F$. For every
short exact sequence of $\mathcal O_X$-modules, they form a morphism from
the algebraic long exact cohomology sequence to the analytic long exact
cohomology sequence. In particular, they commute with connecting maps.
\end{lemma}

\begin{proof}
The derived adjunction unit is a natural transformation. Because
analytification is exact by Lemma
\ref{lemma-analytification-exact-faithful-coherent}, it carries short exact
sequences to short exact sequences. Naturality of the associated
distinguished triangles gives the asserted morphism of long exact
sequences.
\end{proof}


\section{The GAGA comparison theorems}
\label{section-gaga-statements}

\begin{theorem}[Cohomology comparison]
\label{theorem-gaga-cohomology}
\begin{reference}
\cite[\S 3, no. 12, Theorem 1]{GAGA}
\end{reference}
Let $X$ be a projective algebraic variety over $\mathbf C$ and let
$\mathcal F$ be a coherent $\mathcal O_X$-module. For every $q\geq0$, the
map
$$
\epsilon^q_{\mathcal F}:H^q(X,\mathcal F)
\longrightarrow H^q(X^{\mathrm{an}},\mathcal F^{\mathrm{an}})
$$
is an isomorphism. In particular,
$$
\Gamma(X,\mathcal F)\longrightarrow
\Gamma(X^{\mathrm{an}},\mathcal F^{\mathrm{an}})
$$
is an isomorphism.
\end{theorem}

\begin{theorem}[Full faithfulness]
\label{theorem-gaga-fully-faithful}
\begin{reference}
\cite[\S 3, no. 12, Theorem 2]{GAGA}
\end{reference}
Let $X$ be a projective algebraic variety over $\mathbf C$ and let
$\mathcal F,\mathcal G$ be coherent $\mathcal O_X$-modules. Analytification
induces a bijection
$$
\mathop{\rm Hom}\nolimits_{\mathcal O_X}(\mathcal F,\mathcal G)
\longrightarrow
\mathop{\rm Hom}\nolimits_{\mathcal O_{X^{\mathrm{an}}}}
(\mathcal F^{\mathrm{an}},\mathcal G^{\mathrm{an}}).
$$
\end{theorem}

\begin{theorem}[Essential surjectivity]
\label{theorem-gaga-essential-surjectivity}
\begin{reference}
\cite[\S 3, no. 12, Theorem 3]{GAGA}
\end{reference}
Let $X$ be a projective algebraic variety over $\mathbf C$. Every coherent
$\mathcal O_{X^{\mathrm{an}}}$-module $\mathcal M$ is isomorphic to
$\mathcal F^{\mathrm{an}}$ for a coherent $\mathcal O_X$-module
$\mathcal F$. The module $\mathcal F$ is unique up to isomorphism; moreover,
every analytic isomorphism between two such analytifications is induced by
a unique algebraic isomorphism.
\end{theorem}

The proofs of these three theorems occupy the following sections.

\begin{remark}[Projectivity is essential]
\label{remark-gaga-projectivity-essential}
For $X=\mathbf A^1_{\mathbf C}$ and $\mathcal F=\mathcal O_X$, the
degree-zero comparison map is
$$
\mathbf C[z]\longrightarrow
\Gamma(\mathbf C,\mathcal O_{\mathbf C}),
$$
whose target is the ring of entire functions. It is not surjective. Thus the
projectivity hypothesis in the comparison theorems cannot simply be
discarded.
\end{remark}

\begin{remark}[The intermediate inverse-image sheaf]
\label{remark-gaga-comparison-factorization}
The comparison map factors canonically as
$$
H^q(X,\mathcal F)
\longrightarrow H^q(X^{\mathrm{an}},a_X^{-1}\mathcal F)
\longrightarrow H^q(X^{\mathrm{an}},\mathcal F^{\mathrm{an}}).
$$
The first arrow need not be an isomorphism. For example, let $X$ be an
irreducible projective variety whose analytification has a nonzero Betti
number $b_q$ with $q>0$, and put $K=\mathbf C(X)$. The sheaf $K$ is flasque
on the irreducible Zariski space $X$, so $H^q(X,K)=0$. Its inverse image is
the constant sheaf with value $K$ on $X^{\mathrm{an}}$, whose degree-$q$
cohomology is a $K$-vector space of dimension $b_q$. This is the obstruction
recorded after the three theorems in \cite[\S 3, no. 12]{GAGA}.
\end{remark}

% Candidate continuation for GAGA source nos. 13--17.
% Insert after remark-gaga-comparison-factorization and before the
% unique terminal \input{chapters} in the standalone gaga.tex candidate.
% This fragment assumes the analytification functor, its flatness and
% exactness, the comparison maps on cohomology, and the statements of the
% three GAGA theorems have already been introduced.

\section{Proof of the cohomology comparison theorem}
\label{section-proof-gaga-cohomology}

We first record the two analytic inputs used to start and bound the proof.
They are deliberately visible: neither is a formal consequence of the
algebraic cohomology calculation.

\begin{lemma}[Analytic cohomological dimension of projective space]
\label{lemma-projective-analytic-cohomological-dimension}
Let $r \geq 0$ and let $\mathcal A$ be a sheaf of abelian groups on
$\mathbf P^r(\mathbf C)$ with its analytic topology. Then
$$
H^q(\mathbf P^r(\mathbf C),\mathcal A)=0
\quad\text{for }q>2r.
$$
\end{lemma}

\begin{proof}
The underlying space of $\mathbf P^r(\mathbf C)$ is a compact triangulable
real manifold of dimension $2r$. We use the topological cohomological
dimension theorem saying that a paracompact Hausdorff space of covering
dimension $d$ has vanishing sheaf cohomology in degrees greater than $d$.
This covering-dimension theorem is an external topological input; the
current Stacks chapter on Cohomology defines cohomological dimension but
does not prove this comparison with covering dimension.
\end{proof}

\begin{lemma}
\label{lemma-projective-structure-sheaf-comparison}
For every $r \geq 0$ and every $q \geq 0$, the comparison map
$$
H^q(\mathbf P^r_{\mathbf C},\mathcal O)
\longrightarrow
H^q(\mathbf P^r(\mathbf C),\mathcal O^{\mathrm{an}})
$$
is an isomorphism.
\end{lemma}

\begin{proof}
By Cohomology of Schemes, Lemma
\ref{coherent-lemma-cohomology-projective-space-over-ring}, the group on the
left is $\mathbf C$ for $q=0$ and is zero for $q>0$. On the analytic side we
use the Dolbeault theorem together with the standard computation
$$
H^{0,q}_{\overline\partial}(\mathbf P^r(\mathbf C))=
\begin{cases}
\mathbf C & q=0,\\
0 & q>0.
\end{cases}
$$
This Dolbeault--projective-space calculation is the external analytic input
in this lemma; see \cite[\S 2, no. 13]{GAGA}. In degree zero the comparison
map sends a constant algebraic function to the same constant holomorphic
function, and hence is the identity on $\mathbf C$. In positive degree both
sides vanish.
\end{proof}

\begin{lemma}
\label{lemma-projective-twist-comparison}
Let $r \geq 0$ and $n \in \mathbf Z$. For every $q \geq 0$, the comparison
map
$$
H^q(\mathbf P^r_{\mathbf C},\mathcal O(n))
\longrightarrow
H^q(\mathbf P^r(\mathbf C),\mathcal O(n)^{\mathrm{an}})
$$
is an isomorphism.
\end{lemma}

\begin{proof}
We argue by induction on $r$. The assertion for $r=0$ is immediate. Suppose
$r>0$, choose a hyperplane
$i:E\longrightarrow\mathbf P^r_{\mathbf C}$, and identify $E$ with
$\mathbf P^{r-1}_{\mathbf C}$ as in Varieties, Lemma
\ref{varieties-lemma-hyperplane}. Multiplication by an equation of $E$ gives,
for every $n$, a short exact sequence
$$
0\longrightarrow\mathcal O(n-1)\longrightarrow\mathcal O(n)
\longrightarrow i_*\mathcal O_E(n)\longrightarrow0.
$$
Analytification is exact by Lemma
\ref{lemma-analytification-exact-faithful-coherent} and commutes with the
closed-immersion pushforward by Lemma
\ref{lemma-analytification-pushforward-closed-immersion}. Thus this sequence
and its analytification give two long exact cohomology sequences. The
comparison maps form a morphism between them by Lemma
\ref{lemma-analytification-cohomology-map-functorial}.

The induction hypothesis gives the comparison isomorphisms for
$\mathcal O_E(n)$ in all degrees. The four and five lemmas (Homology, Lemmas
\ref{homology-lemma-four-lemma} and \ref{homology-lemma-five-lemma}) now show
that comparison for $\mathcal O(n)$ in every degree is equivalent to
comparison for $\mathcal O(n-1)$ in every degree. The case $n=0$ is Lemma
\ref{lemma-projective-structure-sheaf-comparison}; induction upward and
downward in $n$ proves the result for every integer.
\end{proof}

\begin{proof}[Proof of Theorem \ref{theorem-gaga-cohomology}]
Let $X$ be projective over $\mathbf C$, let $\mathcal F$ be coherent, and
choose a closed immersion $i:X\to\mathbf P^r_{\mathbf C}$. The direct image
$i_*\mathcal F$ is coherent, and Cohomology, Lemma
\ref{cohomology-lemma-cohomology-and-closed-immersions} identifies its
algebraic and analytic cohomology with the corresponding cohomology on $X$
and $X^{\mathrm{an}}$. Lemma
\ref{lemma-analytification-pushforward-closed-immersion} identifies
$(i_*\mathcal F)^{\mathrm{an}}$ with $i_*\mathcal F^{\mathrm{an}}$, compatibly
with the comparison map. We may therefore assume
$X=\mathbf P^r_{\mathbf C}$.

We use descending induction on $q$, simultaneously for all coherent
$\mathcal F$. For $q>2r$, the algebraic group vanishes by Cohomology of
Schemes, Lemma \ref{coherent-lemma-coherent-projective}, and the analytic
group vanishes by Lemma
\ref{lemma-projective-analytic-cohomological-dimension}.

Assume comparison is known in degree $q+1$ for every coherent sheaf. By
Cohomology of Schemes, Lemma
\ref{coherent-lemma-coherent-projective}, there is a short exact sequence
$$
0\longrightarrow\mathcal R\longrightarrow
\mathcal L\longrightarrow\mathcal F\longrightarrow0
$$
with $\mathcal R$ coherent and $\mathcal L$ a finite direct sum of twists
$\mathcal O(n)$. Comparison is an isomorphism for $\mathcal L$ in every
degree by Lemma \ref{lemma-projective-twist-comparison}. In the morphism of
long exact sequences, the four lemma applied to
$$
H^q(\mathcal L)\longrightarrow H^q(\mathcal F)
\longrightarrow H^{q+1}(\mathcal R)\longrightarrow H^{q+1}(\mathcal L)
$$
shows first that the comparison map for $\mathcal F$ in degree $q$ is
surjective. Since this holds for every coherent sheaf, it holds for
$\mathcal R$. Apply the injective half of the four lemma to
$$
H^q(\mathcal R)\longrightarrow H^q(\mathcal L)
\longrightarrow H^q(\mathcal F)\longrightarrow H^{q+1}(\mathcal R).
$$
The comparison map for $\mathcal F$ is injective as well. This completes
the descending induction and the proof.
\end{proof}

\section{Proof of full faithfulness}
\label{section-proof-gaga-fully-faithful}

\begin{lemma}
\label{lemma-analytification-internal-hom}
Let $X$ be a variety over $\mathbf C$ and let $\mathcal F,\mathcal G$ be
coherent $\mathcal O_X$-modules. There is a canonical isomorphism
$$
\SheafHom_{\mathcal O_X}(\mathcal F,\mathcal G)^{\mathrm{an}}
\longrightarrow
\SheafHom_{\mathcal O_{X^{\mathrm{an}}}}
(\mathcal F^{\mathrm{an}},\mathcal G^{\mathrm{an}}).
$$
\end{lemma}

\begin{proof}
Write $a_X:X^{\mathrm{an}}\to X$ for the canonical morphism of locally
ringed spaces. A coherent module on the locally Noetherian scheme $X$ is
finitely presented, and $a_X$ is flat by Theorem
\ref{theorem-analytification-local-faithfully-flat}. The displayed map is
therefore the flat-pullback internal-Hom isomorphism of Modules, Lemma
\ref{modules-lemma-pullback-internal-hom}. Equivalently, on a stalk it is
the flat base-change isomorphism of More on Algebra, Lemma
\ref{more-algebra-lemma-pseudo-coherence-and-base-change-ext}. Coherence of
the source follows from Cohomology of Schemes, Lemma
\ref{coherent-lemma-tensor-hom-coherent}; coherence of the target is also the
$q=0$ case of Cohomology, Lemma
\ref{cohomology-lemma-coherent-analytic-sheaf-ext}.
\end{proof}

\begin{proof}[Proof of Theorem \ref{theorem-gaga-fully-faithful}]
Put $\mathcal A=\SheafHom_{\mathcal O_X}(\mathcal F,\mathcal G)$. Then
$$
\Hom_{\mathcal O_X}(\mathcal F,\mathcal G)
=H^0(X,\mathcal A).
$$
Theorem \ref{theorem-gaga-cohomology} in degree zero and Lemma
\ref{lemma-analytification-internal-hom} give canonical isomorphisms
$$
H^0(X,\mathcal A)
\longrightarrow H^0(X^{\mathrm{an}},\mathcal A^{\mathrm{an}})
\longrightarrow
\Hom_{\mathcal O_{X^{\mathrm{an}}}}
(\mathcal F^{\mathrm{an}},\mathcal G^{\mathrm{an}}).
$$
Their composite sends $f$ to $f^{\mathrm{an}}$ by construction of the
cohomology comparison map. Hence analytification is fully faithful.
\end{proof}

\section{Essential surjectivity: preliminaries and reduction}
\label{section-gaga-essential-surjectivity-preliminaries}

\begin{lemma}
\label{lemma-coherent-analytic-closed-immersion}
Let $i:Y\to X$ be a closed immersion of complex analytic spaces whose
structure sheaves and defining ideal are coherent. Then $i_*$ is exact and
fully faithful on modules, carries coherent $\mathcal O_Y$-modules to
coherent $\mathcal O_X$-modules, and its essential image consists of the
coherent modules annihilated by the defining ideal. Moreover,
$$
H^q(Y,\mathcal M)=H^q(X,i_*\mathcal M)
$$
for every $q\geq0$.
\end{lemma}

\begin{proof}
Exactness, full faithfulness, and the description of the essential image are
Modules, Lemma \ref{modules-lemma-i-star-equivalence}. To check coherence,
first regard $i_*\mathcal M$ as a module over
$i_*\mathcal O_Y=\mathcal O_X/\mathcal I$. Modules, Lemma
\ref{modules-lemma-i-star-reflects-finite-type}, applied with this quotient
structure sheaf, identifies coherence on $Y$ with coherence over
$i_*\mathcal O_Y$. Modules, Lemma
\ref{modules-lemma-coherent-change-rings} then identifies this with
coherence over $\mathcal O_X$. The cohomology equality is Cohomology, Lemma
\ref{cohomology-lemma-cohomology-and-closed-immersions}.
\end{proof}

\begin{lemma}
\label{lemma-gaga-essential-surjectivity-projective-reduction}
Uniqueness in Theorem \ref{theorem-gaga-essential-surjectivity} follows from
Theorem \ref{theorem-gaga-fully-faithful}. For existence, it is enough to
prove the theorem when $X=\mathbf P^r_{\mathbf C}$.
\end{lemma}

\begin{proof}
If $\mathcal F^{\mathrm{an}}\cong\mathcal G^{\mathrm{an}}$, full faithfulness
lifts the isomorphism and its inverse uniquely to algebraic maps. Faithfulness
shows that their composites are the identity, proving uniqueness.

For the reduction, let $i:Y\to\mathbf P^r_{\mathbf C}$ be a closed immersion
and let $\mathcal M$ be coherent on $Y^{\mathrm{an}}$. By Lemma
\ref{lemma-coherent-analytic-closed-immersion}, $i_*\mathcal M$ is coherent
on $\mathbf P^r(\mathbf C)$. Assuming essential surjectivity on projective
space, choose a coherent algebraic module $\mathcal G$ with
$\mathcal G^{\mathrm{an}}\cong i_*\mathcal M$.

Let $\mathcal I$ be the ideal of $Y$. Compatibility of analytification with
ideals, tensor products, and images identifies
$(\mathcal I\mathcal G)^{\mathrm{an}}$ with
$\mathcal I^{\mathrm{an}}\mathcal G^{\mathrm{an}}$, which is zero because
$i_*\mathcal M$ is annihilated by $\mathcal I^{\mathrm{an}}$. Exactness and
faithfulness in Lemma
\ref{lemma-analytification-exact-faithful-coherent} therefore give
$\mathcal I\mathcal G=0$. Morphisms, Lemma
\ref{morphisms-lemma-i-star-equivalence} and Cohomology of Schemes, Lemma
\ref{coherent-lemma-i-star-equivalence} give a coherent algebraic
$\mathcal F$ on $Y$ with $i_*\mathcal F=\mathcal G$. Finally,
Lemma \ref{lemma-analytification-pushforward-closed-immersion} gives
$$
i_*\mathcal F^{\mathrm{an}}
\cong(i_*\mathcal F)^{\mathrm{an}}
\cong\mathcal G^{\mathrm{an}}
\cong i_*\mathcal M.
$$
Full faithfulness of analytic closed-immersion pushforward yields
$\mathcal F^{\mathrm{an}}\cong\mathcal M$.
\end{proof}

\section{Analytic twists and global generation}
\label{section-gaga-analytic-twists}

We now fix $X=\mathbf P^r_{\mathbf C}$ and argue by induction on $r$ for
essential surjectivity. The case $r=0$ is the equivalence between finite
dimensional complex vector spaces and coherent modules on a point.

\begin{definition}
\label{definition-coherent-analytic-twist}
For a coherent analytic module $\mathcal M$ on $X^{\mathrm{an}}$ and
$n\in\mathbf Z$, set
$$
\mathcal M(n)=
\mathcal M\otimes_{\mathcal O_{X^{\mathrm{an}}}}
\mathcal O_X(n)^{\mathrm{an}}.
$$
This is coherent, and the tensor compatibility of pullback gives the
canonical identity
$$
\mathcal F^{\mathrm{an}}(n)\cong\mathcal F(n)^{\mathrm{an}}
$$
for every coherent algebraic $\mathcal F$.
\end{definition}

\begin{lemma}
\label{lemma-hyperplane-analytic-serre-vanishing}
Let $E\cong\mathbf P^{r-1}_{\mathbf C}$ be a hyperplane and let $\mathcal A$
be a coherent analytic module on $E^{\mathrm{an}}$. Then
$$
H^q(E^{\mathrm{an}},\mathcal A(n))=0
$$
for every $q>0$ and all sufficiently large $n$.
\end{lemma}

\begin{proof}
By the induction hypothesis for essential surjectivity there is a coherent
algebraic $\mathcal F$ on $E$ with
$\mathcal A\cong\mathcal F^{\mathrm{an}}$. Definition
\ref{definition-coherent-analytic-twist} and Theorem
\ref{theorem-gaga-cohomology} identify the displayed group with
$H^q(E,\mathcal F(n))$. The latter vanishes for $q>0$ and $n$ sufficiently
large by Cohomology of Schemes, Lemma
\ref{coherent-lemma-coherent-projective}. There are only finitely many
possibly nonzero degrees, so one lower bound on $n$ works for all $q>0$.
\end{proof}

\begin{lemma}
\label{lemma-projective-analytic-global-generation}
Let $\mathcal M$ be a coherent analytic module on
$X^{\mathrm{an}}=\mathbf P^r(\mathbf C)$. There is an integer $n_0$ such
that $\mathcal M(n)$ is generated by global sections for every $n\geq n_0$.
\end{lemma}

\begin{proof}
We first make two elementary observations. If global sections generate
$\mathcal M(n)$ at $x$, choose a homogeneous coordinate $T_k$ nonzero at
$x$. Multiplication by $T_k^{m-n}$ sends global sections of $\mathcal M(n)$
to global sections of $\mathcal M(m)$ and is an isomorphism on the stalk at
$x$; generation at $x$ therefore persists for every $m\geq n$. Generation
at a point is open: choose finitely many global sections generating the
stalk, take their evaluation map from a finite free module, and use
coherence of its cokernel to see that it vanishes near the point. Since
$X^{\mathrm{an}}$ is compact, it is consequently enough to find, for every
$x$, one twist generated at $x$.

Fix $x$ and choose a hyperplane $i:E\to X$ through it, cut out by a linear
form $t$. Tensoring
$0\to\mathcal O_X(-1)^{\mathrm{an}}\xrightarrow{t}
\mathcal O_{X^{\mathrm{an}}}\to i_*\mathcal O_{E^{\mathrm{an}}}\to0$
with $\mathcal M$ gives an exact sequence
$$
0\longrightarrow\mathcal C\longrightarrow\mathcal M(-1)
\xrightarrow{t}\mathcal M\longrightarrow\mathcal B\longrightarrow0,
$$
where $\mathcal C$ and $\mathcal B$ are coherent. Both are annihilated by
$t$, so Lemma \ref{lemma-coherent-analytic-closed-immersion} regards them
as coherent modules on $E^{\mathrm{an}}$.

After twisting by $n$, let $\mathcal L_n$ be the image of
$\mathcal M(n-1)\to\mathcal M(n)$. We have short exact sequences
$$
0\to\mathcal C(n)\to\mathcal M(n-1)\to\mathcal L_n\to0,
\qquad
0\to\mathcal L_n\to\mathcal M(n)\to\mathcal B(n)\to0.
$$
Lemma \ref{lemma-hyperplane-analytic-serre-vanishing} gives, for all large
$n$,
$$
H^2(X^{\mathrm{an}},\mathcal C(n))=0,
\qquad
H^1(X^{\mathrm{an}},\mathcal B(n))=0.
$$
The two long exact sequences then give
$$
\dim H^1(\mathcal M(n-1))
\geq\dim H^1(\mathcal L_n)
\geq\dim H^1(\mathcal M(n)).
$$
All these dimensions are finite by Cohomology, Lemma
\ref{cohomology-lemma-compact-complex-analytic-coherent-cohomology-finite-dimensional}.
Thus the nonnegative integers $\dim H^1(\mathcal M(n))$ eventually
stabilize. For $n$ in the stable range, the surjection
$H^1(\mathcal L_n)\to H^1(\mathcal M(n))$ is an isomorphism, and the long
exact sequence yields a surjection
$$
H^0(X^{\mathrm{an}},\mathcal M(n))
\longrightarrow H^0(X^{\mathrm{an}},\mathcal B(n)).
$$

By induction on $r$, the module on $E^{\mathrm{an}}$ corresponding to
$\mathcal B$ is the analytification of a coherent algebraic module
$\mathcal G$ on $E$. For all sufficiently large $n$, the module
$\mathcal G(n)$ is generated by global sections by Properties, Proposition
\ref{properties-proposition-characterize-ample}. Theorem
\ref{theorem-gaga-cohomology} in degree zero says the same for
$\mathcal B(n)$, and the preceding surjection lifts those generators to
global sections of $\mathcal M(n)$.

At $x$, put $A=\mathcal O_{X^{\mathrm{an}},x}$,
$M=\mathcal M(n)_x$, and let $N\subset M$ be the submodule generated by
global sections. If $\mathfrak p=(t)\subset A$, then
$\mathcal B(n)_x=M/\mathfrak pM$. The lifted sections generate this
quotient, so $M=N+\mathfrak pM$. Since $\mathfrak p$ lies in the maximal
ideal of the local ring $A$, Nakayama's lemma (Algebra, Lemma
\ref{algebra-lemma-NAK}) gives $M=N$. This supplies the required twist at
$x$; the two opening observations and compactness supply one $n_0$ valid at
every point and every $n\geq n_0$.
\end{proof}

\section{Completion of essential surjectivity}
\label{section-completion-gaga-essential-surjectivity}

\begin{proof}[Completion of the proof of Theorem
\ref{theorem-gaga-essential-surjectivity}]
It remains, by Lemma
\ref{lemma-gaga-essential-surjectivity-projective-reduction}, to treat
$X=\mathbf P^r_{\mathbf C}$. Let $\mathcal M$ be coherent on
$X^{\mathrm{an}}$. Lemma
\ref{lemma-projective-analytic-global-generation} and compactness give a
finite set of global generators after a twist. Untwisting produces a
surjection
$$
\mathcal L_0^{\mathrm{an}}\longrightarrow\mathcal M,
\qquad
\mathcal L_0=\mathcal O_X(-n)^{\oplus p}.
$$
Its kernel $\mathcal R$ is coherent by Modules, Lemma
\ref{modules-lemma-coherent-abelian}. Apply the same global-generation
argument to $\mathcal R$ to obtain a finite direct sum of twists
$\mathcal L_1$ and a surjection
$\mathcal L_1^{\mathrm{an}}\to\mathcal R$. Composing with the inclusion of
$\mathcal R$ gives an exact sequence
$$
\mathcal L_1^{\mathrm{an}}\xrightarrow{g}
\mathcal L_0^{\mathrm{an}}\longrightarrow\mathcal M\longrightarrow0.
$$

Theorem \ref{theorem-gaga-fully-faithful} supplies a unique algebraic map
$f:\mathcal L_1\to\mathcal L_0$ with $f^{\mathrm{an}}=g$. Put
$\mathcal F=\mathop{\rm Coker}(f)$. This is coherent, and exactness of
analytification gives
$$
\mathcal F^{\mathrm{an}}
=\mathop{\rm Coker}(f^{\mathrm{an}})
=\mathop{\rm Coker}(g)
\cong\mathcal M.
$$
Together with the reduction and uniqueness already proved, this completes
Theorem \ref{theorem-gaga-essential-surjectivity}.
\end{proof}

\section{Applications of the comparison theorems}
\label{section-gaga-applications}

\subsection{Conjugation and Betti numbers}
\label{subsection-gaga-betti-numbers}

Let $\sigma\in\mathop{\rm Aut}(\mathbf C)$. For a scheme $X$ over
$\mathbf C$, write
$$
X^\sigma=X\times_{\mathop{\rm Spec}(\mathbf C),\sigma}
\mathop{\rm Spec}(\mathbf C).
$$

\begin{proposition}
\label{proposition-conjugate-smooth-projective-betti-numbers}
\begin{reference}
\cite[\S 4, no. 18, Proposition 12]{GAGA}
\end{reference}
Let $X$ be a smooth projective algebraic variety over $\mathbf C$. Then $X^{\rm an}$
and $(X^\sigma)^{\rm an}$ have the same Betti numbers.
\end{proposition}

\begin{proof}
For $p,q\geq 0$, flat base change gives a $\sigma$-semilinear
isomorphism
$$
H^q(X,\Omega^p_{X/\mathbf C})
\longrightarrow
H^q(X^\sigma,\Omega^p_{X^\sigma/\mathbf C});
$$
see Cohomology of Schemes, Lemma
\ref{coherent-lemma-flat-base-change-cohomology}. In particular, the two
vector spaces have the same dimension. For a smooth variety, local
coordinates give the canonical comparison
$$
(\Omega^p_{X/\mathbf C})^{\rm an}
\cong \Omega^p_{X^{\rm an}},
$$
where the module on the right is the module of holomorphic $p$-forms. Thus
the cohomology comparison theorem \ref{theorem-gaga-cohomology} identifies
these dimensions with
$$
h^{p,q}(X^{\rm an})
\quad\hbox{and}\quad
h^{p,q}((X^\sigma)^{\rm an}),
$$
respectively, where
$$
h^{p,q}(M)=\dim_{\mathbf C}H^q(M,\Omega^p_M).
$$
Finally, a smooth projective analytification is compact K\"ahler. Hodge
decomposition (using the Dolbeault resolution and the K\"ahler identities),
which is an external analytic input here, gives
$$
b_n(X^{\rm an})=\sum_{p+q=n}h^{p,q}(X^{\rm an})
$$
and the analogous formula for $X^\sigma$. The assertion follows.
\end{proof}

\begin{proposition}[Corollary]
\label{corollary-number-field-embeddings-betti-numbers}
\begin{reference}
\cite[\S 4, no. 18, corollary to Proposition 12]{GAGA}
\end{reference}
Let $V$ be a smooth projective variety over a field $K$ algebraic over
$\mathbf Q$. The Betti
numbers of $V\times_{K,\iota}\mathbf C$ are independent of the embedding
$\iota:K\longrightarrow\mathbf C$.
\end{proposition}

\begin{proof}
The standard extension theorem for field embeddings shows that two
embeddings of $K$ in $\mathbf C$ differ by an automorphism of $\mathbf C$.
Apply Proposition
\ref{proposition-conjugate-smooth-projective-betti-numbers}.
\end{proof}

\begin{remark}[Historical scope]
\label{remark-conjugate-varieties-historical}
Serre observed in \cite[\S 4, no. 18]{GAGA} that a conjugate pair need not be
analytically isomorphic (elliptic curves already show this), and asked there
whether such a pair must be homeomorphic. This historical question is not an
input to Proposition
\ref{proposition-conjugate-smooth-projective-betti-numbers}.
\end{remark}


\subsection{Chow's theorem and algebraicity}
\label{subsection-gaga-chow}

\begin{theorem}[Chow]
\label{theorem-chow-closed-analytic-projective}
\begin{reference}
\cite[\S 4, no. 19, Proposition 13]{GAGA}
\end{reference}
Let $X$ be a projective algebraic variety over $\mathbf C$. Every closed reduced
analytic subspace $Y\subset X^{\rm an}$ is the analytification of a unique
closed reduced subscheme of $X$.
\end{theorem}

\begin{proof}
The vanishing ideal $\mathcal J_Y\subset\mathcal O_{X^{\rm an}}$ is
coherent by Theorem \ref{theorem-oka-cartan-coherence}. By Theorem
\ref{theorem-gaga-essential-surjectivity}, there is a coherent
$\mathcal O_X$-module $\mathcal F$ with
$\mathcal F^{\rm an}\cong\mathcal J_Y$. Compose this isomorphism with the
inclusion into $\mathcal O_{X^{\rm an}}$. The full-faithfulness theorem
\ref{theorem-gaga-fully-faithful} algebraizes the composite to an
$\mathcal O_X$-linear map
$$
u:\mathcal F\longrightarrow\mathcal O_X.
$$
Its image $\mathcal I$ is a coherent ideal. Exactness and faithfulness of
analytification identify $\mathcal I^{\rm an}$ with $\mathcal J_Y$. Hence
the ideal $\mathcal I$ is radical. Indeed, at every closed point
$x\in X(\mathbf C)$, if $a^m\in\mathcal I_x$, then the image of $a$ in
$\mathcal O_{X^{\rm an},x}$ belongs to the radical ideal
$\mathcal J_{Y,x}=\mathcal I_x\mathcal O_{X^{\rm an},x}$. Faithful
flatness in Theorem \ref{theorem-analytification-local-faithfully-flat}
contracts this ideal back to $\mathcal I_x$. Since $X$ is Noetherian and
Jacobson, radicality at all closed stalks implies that $\mathcal I$ is
radical. The closed reduced subscheme defined by $\mathcal I$ therefore has
analytic ideal $\mathcal J_Y$, and hence has analytification $Y$.

If two reduced closed subschemes have analytification $Y$, their ideal
sheaves have equal extensions at every analytic local ring. Faithful
flatness contracts these extensions to equal stalks at every closed point;
coherence and the Jacobson property then identify the ideal sheaves. This
proves uniqueness.
\end{proof}

\begin{proposition}
\label{proposition-compact-analytic-subspace-algebraic}
\begin{reference}
\cite[\S 4, no. 19, Proposition 14]{GAGA}
\end{reference}
Let $X$ be an algebraic variety over $\mathbf C$. Every compact
closed reduced analytic subspace $Y\subset X^{\rm an}$ is algebraic: there
is a unique closed reduced subscheme $Z\subset X$ with
$Z^{\rm an}=Y$.
\end{proposition}

\begin{proof}
Nagata compactification and Chow's lemma (More on Flatness, Theorem
\ref{flat-theorem-nagata} and Limits of Schemes, Lemma
\ref{limits-lemma-chow-finite-type}) give an open immersion
$X\subset\overline X$ with $\overline X$ proper and a proper surjection
$p:X'\to\overline X$ together with an immersion
$X'\to\mathbf P^n_{\mathbf C}$. Replacing $X'$ by its reduction preserves
these properties. The resulting variety $X'$ is proper over $\mathbf C$,
so its immersion is closed and $X'$ is projective. Since $Y$ is compact,
it is closed in the Hausdorff space $\overline X^{\rm an}$. Its analytic structure on
$X^{\rm an}$ glues with the empty closed analytic subspace on
$\overline X^{\rm an}\setminus Y$, and therefore makes $Y$ a closed reduced
analytic subspace of $\overline X^{\rm an}$.

The reduced inverse image
$$
Y'=\left((X')^{\rm an}\times_{\overline X^{\rm an}}Y\right)_{\rm red}
$$
is a closed reduced analytic subspace of $(X')^{\rm an}$. Theorem
\ref{theorem-chow-closed-analytic-projective} algebraizes it to a closed
reduced subscheme $Z'\subset X'$. Let $Z\subset\overline X$ be the reduced
Zariski closure of $p(Z')$. The analytic map sends $(Z')^{\rm an}=Y'$ onto
$Y$. Since $Y$ is closed, Proposition \ref{proposition-image-closure}
shows that the classical point set of $Z^{\rm an}$ is exactly $Y$.
Reduced closed analytic subspaces are determined by their underlying
analytic sets, so $Z^{\rm an}=Y$. All closed points of $Z$ lie in $X$;
the Jacobson property therefore shows that $Z\subset X$.

For uniqueness, if $Z_1,Z_2\subset X$ have analytification $Y$, their ideal
stalks have equal extensions to every $\mathcal O_{X^{\rm an},x}$.
Faithful flatness in Theorem
\ref{theorem-analytification-local-faithfully-flat} contracts them to equal
stalks at every closed point. Coherence and the Jacobson property give
$Z_1=Z_2$.
\end{proof}

\begin{proposition}
\label{proposition-holomorphic-map-from-proper-is-regular}
\begin{reference}
\cite[\S 4, no. 19, Proposition 15]{GAGA}
\end{reference}
Let $X$ be a proper algebraic variety over $\mathbf C$ and let $Y$ be an
algebraic variety over $\mathbf C$. Every holomorphic map
$$
f:X^{\rm an}\longrightarrow Y^{\rm an}
$$
is the analytification of a unique morphism $X\to Y$.
\end{proposition}

\begin{proof}
The graph $\Gamma_f$ is a compact closed analytic subspace of
$(X\times_{\mathbf C}Y)^{\rm an}$. Proposition
\ref{proposition-compact-analytic-subspace-algebraic} makes it the
analytification of a closed reduced subvariety of $X\times_{\mathbf C}Y$.
Proposition \ref{proposition-algebraic-graph-holomorphic-regular} therefore
defines the required algebraic map. For uniqueness, two candidate
maps have closed reduced graphs with the same classical points; the graphs,
and hence the maps, are equal.
\end{proof}

\begin{proposition}[Corollary]
\label{corollary-compact-analytic-space-unique-algebraization}
\begin{reference}
\cite[\S 4, no. 19, corollary to Proposition 15]{GAGA}
\end{reference}
A compact reduced analytic space admits at most one algebraic-variety
structure: any analytic isomorphism between two algebraizations is the
analytification of a unique algebraic isomorphism.
\end{proposition}

\begin{proof}
Every algebraization is proper by Proposition
\ref{proposition-proper-if-and-only-if-compact-analytification}. Apply
Proposition \ref{proposition-holomorphic-map-from-proper-is-regular} to the
analytic isomorphism and to its inverse.
\end{proof}


\subsection{Algebraic and analytic principal bundles}
\label{subsection-gaga-principal-bundles}

\begin{definition}
\label{definition-principal-bundle-analytification-map}
\begin{reference}
\cite[\S 4, no. 20]{GAGA}
\end{reference}
Let $G$ be an algebraic group over $\mathbf C$ and let $X$ be an algebraic
variety over $\mathbf C$. Write
$$
H^1_{\rm Zar}(X,G)
$$
for the pointed set of isomorphism classes of Zariski locally trivial
principal $G$-bundles and
$$
H^1_{\rm an}(X^{\rm an},G^{\rm an})
$$
for the pointed set of isomorphism classes of holomorphically locally
trivial principal $G^{\rm an}$-bundles. Analytification defines a map
$$
\varepsilon_G:H^1_{\rm Zar}(X,G)
\longrightarrow H^1_{\rm an}(X^{\rm an},G^{\rm an}).
$$
This is the torsor interpretation of first cohomology from Cohomology on
Sites, Section \ref{sites-cohomology-section-h1-torsors}.
\end{definition}

\begin{proposition}
\label{proposition-principal-bundle-analytification-injective}
\begin{reference}
\cite[\S 4, no. 20, Proposition 16]{GAGA}
\end{reference}
If $X$ is a proper algebraic variety over $\mathbf C$, then $\varepsilon_G$
is injective. More
precisely, every analytic isomorphism between the analytifications of two
algebraic principal $G$-bundles is algebraic.
\end{proposition}

\begin{proof}
For two algebraic principal $G$-bundles $P$ and $P'$ on $X$, the associated
bundle
$$
\mathop{\rm Isom}\nolimits_G(P,P')
$$
is an algebraic fibre bundle over $X$, and its sections are precisely the
$G$-equivariant isomorphisms $P\to P'$. An analytic isomorphism gives a
holomorphic section of its analytification. Proposition
\ref{proposition-holomorphic-map-from-proper-is-regular} makes the section
algebraic.
\end{proof}

\begin{proposition}
\label{proposition-additive-principal-bundles-gaga}
\begin{reference}
\cite[\S 4, no. 20, Proposition 17]{GAGA}
\end{reference}
If $X$ is a projective algebraic variety over $\mathbf C$ and
$G=\mathbf G_a$, then
$\varepsilon_G$ is bijective.
\end{proposition}

\begin{proof}
Principal $\mathbf G_a$-bundles are torsors under the additive sheaf
$\mathcal O_X$, and hence are classified by $H^1(X,\mathcal O_X)$; see
Cohomology, Lemma \ref{cohomology-lemma-torsors-h1}. The analytic statement
uses $H^1(X^{\rm an},\mathcal O_{X^{\rm an}})$. The comparison isomorphism
of Theorem \ref{theorem-gaga-cohomology} identifies these groups.
\end{proof}

\begin{proposition}
\label{proposition-general-linear-principal-bundles-gaga}
\begin{reference}
\cite[\S 4, no. 20, Proposition 18]{GAGA}
\end{reference}
If $X$ is a projective algebraic variety over $\mathbf C$ and
$G=\mathop{\rm GL}_n$, then $\varepsilon_G$ is bijective.
\end{proposition}

\begin{proof}
Principal $\mathop{\rm GL}_n$-bundles are equivalent to rank $n$ vector
bundles. Let $\mathcal E$ be the locally free coherent analytic module
associated with an analytic vector bundle. By Theorem
\ref{theorem-gaga-essential-surjectivity}, it is isomorphic to
$\mathcal F^{\rm an}$ for a coherent $\mathcal O_X$-module $\mathcal F$.
For every closed point $x\in X(\mathbf C)$, the faithfully flat local map
of Theorem \ref{theorem-analytification-local-faithfully-flat} makes
$$
\mathcal F_x\otimes_{\mathcal O_{X,x}}
\mathcal O_{X^{\rm an},x}
$$
free of rank $n$. Faithfully flat descent for finite projectivity (Algebra,
Proposition \ref{algebra-proposition-ffdescent-finite-projectivity}) shows
that $\mathcal F_x$ is finite projective. Algebra, Theorem
\ref{algebra-theorem-projective-free-over-local-ring} makes it free, and
faithful base change fixes its rank as $n$. Since $X$ is Jacobson and
$\mathcal F$ is coherent, local freeness at every closed point makes
$\mathcal F$ locally free of rank $n$ everywhere. Its frame bundle
algebraizes the given analytic principal bundle. Injectivity is Proposition
\ref{proposition-principal-bundle-analytification-injective}.
\end{proof}

\begin{remark}[Line bundles and the historical formulation]
\label{remark-line-bundles-gaga}
For $n=1$, Proposition
\ref{proposition-general-linear-principal-bundles-gaga} is the isomorphism
$$
\mathop{\rm Pic}(X)\longrightarrow\mathop{\rm Pic}(X^{\rm an}).
$$
Serre's remarks after Proposition 18 in \cite[\S 4, no. 20]{GAGA} express
this in terms of divisor classes on a normal variety and record the earlier
smooth result of Kodaira--Spencer. Those historical attributions are not
additional hypotheses here. Serre also notes that Proposition 18 extends
Kodaira's Theorems 7 and 8 to arbitrary projective varieties, possibly with
singularities, without pursuing those applications.
\end{remark}

\begin{proposition}[Algebraizing a reduction of structure group]
\label{proposition-algebraize-reduction-structure-group}
\begin{reference}
\cite[\S 4, no. 20, Proposition 19]{GAGA}
\end{reference}
Let $H\subset G$ be algebraic groups over $\mathbf C$. Assume that the
quotient $G/H$ exists as an algebraic variety and that
$G\to G/H$ has a rational section. Let $X$ be a proper algebraic variety
over $\mathbf C$ and let $P$ be an analytic principal $H^{\rm an}$-bundle
on $X^{\rm an}$. Then
$P$ is algebraizable if and only if the extended bundle
$$
P\times^H G
$$
is algebraizable.
\end{proposition}

\begin{proof}
A rational section over a nonempty open of $G/H$, and the translates of
its domain by $G$, show that $G\to G/H$ is Zariski locally trivial as an
$H$-bundle.
Necessity is therefore immediate by extension of structure group.

Conversely, choose an algebraic principal $G$-bundle $Q$ and an analytic
isomorphism $Q^{\rm an}\cong P\times^H G$. Form the algebraic associated
bundle
$$
E=Q\times^G(G/H).
$$
The analytic reduction $P$ determines a holomorphic section
$s:X^{\rm an}\to E^{\rm an}$. By Proposition
\ref{proposition-holomorphic-map-from-proper-is-regular}, this section is
algebraic. Pulling the principal $H$-bundle $Q\to E$ back along $s$ gives an
algebraic principal $H$-bundle whose analytification is $P$.
\end{proof}

\begin{definition}
\label{definition-rational-section-condition}
\begin{reference}
\cite[\S 4, no. 20, condition (R)]{GAGA}
\end{reference}
A closed algebraic subgroup $G\subset\mathop{\rm GL}_n$ satisfies
\textup{(R)} if the quotient exists and the projection
$$
\mathop{\rm GL}_n/G\longleftarrow\mathop{\rm GL}_n
$$
has a rational section.
\end{definition}

\begin{proposition}
\label{proposition-rational-section-principal-bundles-gaga}
\begin{reference}
\cite[\S 4, no. 20, Proposition 20]{GAGA}
\end{reference}
Let $G\subset\mathop{\rm GL}_n$ satisfy \textup{(R)}. For every projective
algebraic variety $X$ over $\mathbf C$, the map
$$
\varepsilon_G:H^1_{\rm Zar}(X,G)
\longrightarrow H^1_{\rm an}(X^{\rm an},G^{\rm an})
$$
is bijective.
\end{proposition}

\begin{proof}
Given an analytic principal $G$-bundle $P$, Proposition
\ref{proposition-general-linear-principal-bundles-gaga} algebraizes
$P\times^G\mathop{\rm GL}_n$. Proposition
\ref{proposition-algebraize-reduction-structure-group}, with
$H=G$ and the ambient group $\mathop{\rm GL}_n$, then algebraizes $P$.
Injectivity is Proposition
\ref{proposition-principal-bundle-analytification-injective}.
\end{proof}

\begin{remark}[Examples and historical boundary]
\label{remark-rational-section-condition-examples}
Condition \textup{(R)} holds for solvable subgroups by Rosenlicht's theorem,
for $\mathop{\rm SL}_n$ via the determinant section, and for
$\mathop{\rm Sp}_{2m}$ by rational symplectic Gram--Schmidt. Serre also
records in \cite[\S 4, no. 20]{GAGA} the then-open simply connected
semisimple conjecture and the failure of \textup{(R)} for the unimodular
orthogonal group in dimensions at least $3$; in the latter case he leaves
bijectivity of $\varepsilon_G$ open. These historical observations are not
used in Proposition
\ref{proposition-rational-section-principal-bundles-gaga}.
\end{remark}

\input{chapters}

\bibliography{my}
\bibliographystyle{amsalpha}

\end{document}
